% critique.tex — f167r, vv. 380–422 : commentaire critique
% Articles textuels tirés du commentaire de l'auteur (latex/Aristophane-Thesmo Harmattan 2023 vv 1-641.tex),
% texte et notes conservés ; répartition détaillée dans repartition.md.
\cc{393}{τὰς οἰνοπίπας,}{Encore une fois, c'est le scholiaste, disposant de meilleurs manuscrits, qui nous permet de revenir à la leçon originale : après avoir évoqué Symmaque, il rapproche non sans raison \grec{οἰνοπίπης} de \grec{παιδοπίπης} (on aurait pu rajouter : \grec{γυναικοπίπης}, \grec{παρθενοπίπης} ou encore le \grec{πυροπίπης} / πυρροπίπης des \italique{Cavaliers}, v. 407) ; c'est-à-dire que le suffixe ne renverrait pas aux verbes \grec{πιπίζειν} ou \grec{πιπίσκειν} mais -\grec{οπίπης} serait de la famille de \grec{ὀπιπεύειν}, \italique{guigner}, \italique{guetter}, \italique{épier} (cf. aussi Hésychius (ο 1006) : \grec{ὀπῖπα· ἐξαπατᾷ, ἀπατεών, ἢ ἀπατῶν}).}
\cc{400}{τοὺς ἄνδρας ἡμῶν· ὥστ᾽ ἐάν τις πλέκῃ.}{Le vers 400 pèche contre le mètre. Willems y supplée en lisant <ξυμ>πλέκῃ ; il s'en justifie dans la \nieme{11} note, reproduite à la fin.}
\cc{418–20}{ἃ δ᾽ ἦν ἡμῖν πρὸ τοῦ / αὐταῖς ταμιεύεσθαι προαιρούσαις λαβεῖν / ἄλφιτον ἔλαιον οἶνον.}{Je renvoie à la courte mais incisive note (la \nieme{12}) de Willems reproduite à la fin : il y fait un sort à la « déplorable conjecture de Scaliger » (qui avait substitué λαθεῖν à λαβεῖν au vers 419, donnant au passage un sens qui fait contre-sens).}
